% from ega2-1-fr.tex
\oldpage[II]{25}
on a $\overline{x}=0$. D'autre part, si $x\in S_{nd}$ est tel que
$x=(f-1)y$ avec
$y=y_{hd}+y_{(h+1)d}+\cdots+y_{kd}$ avec $y_{jd}\in S_{jd}$ et
$y_{hd}\neq 0$, on a nécessairement $h=n$ et $x=-y_{hd}$, ainsi que les
relations $y_{(j+1)d}=fy_{jd}$ pour $h\leq j\leq k-1$, $fy_{kd}=0$, ce qui
donne finalement $f^{k-n+1}x=0$~; à toute classe $\overline{x}$ mod.
$(f-1)S^{(d)}$ d'un élément $x\in S_{nd}$, on peut donc faire correspondre
l'élément $x/f^n$ de $S_{(f)}$, car la remarque qui précède montre que cette
application est bien définie. Il est immédiat que les deux applications ainsi
définies sont des homomorphismes d'anneaux, réciproques l'un de l'autre. On
procède exactement de même pour $M$.
\end{proof}

\begin{corollary}[2.2.6]
\label{II.2.2.6-fr}
Si $S$ est noethérien, il en est de même de $S_{(f)}$ pour tout $f$ homogène
de degré $>0$.
\end{corollary}

\begin{proof}
Cela résulte aussitôt de (\hyperref[II.2.1.7-fr]{2.1.7}) et
(\hyperref[II.2.2.5-fr]{2.2.5}).
\end{proof}

\begin{env}[2.2.7]
\label{II.2.2.7-fr}
Soit $T$ une partie multiplicative de $S_+$ formée d'éléments
\emph{homogènes}~; $T_0=T\cup\{1\}$ est alors une partie multiplicative de
$S$~; comme les éléments de $T_0$ sont homogènes, l'anneau $T_0^{-1}S$ est
encore gradué de façon évidente~; on désignera par $S_{(T)}$ le sous-anneau
de $T_0^{-1}S$ formé des éléments d'ordre $0$, c'est-à-dire des éléments de
la forme $x/h$, où $h\in T$ et $x$ est homogène de degré égal à celui de
$h$. On sait (\hyperref[0.1.4.5-fr]{0, 1.4.5}) que $T_0^{-1}S$ s'identifie
canoniquement à la limite inductive des anneaux $S_f$, où $f$ parcourt $T$
(pour les homomorphismes canoniques $S_f\to S_{fg}$)~; comme cette
identification respecte les degrés, elle identifie $S_{(T)}$ à la
\emph{limite inductive} des $S_{(f)}$ pour $f\in T$. Pour tout $S$-module
gradué $M$, on définit de même le module $M_{(T)}$ (sur l'anneau $S_{(T)}$)
formé des éléments de degré $0$ de $T_0^{-1}M$, et on voit que ce module est
limite inductive des $M_{(f)}$ pour $f\in T$.

Si $\mathfrak{p}$ est un idéal premier gradué de $S_+$, on désignera par
$S_{(\mathfrak{p})}$ et $M_{(\mathfrak{p})}$ l'anneau $S_{(T)}$ et le
module $M_{(T)}$ respectivement, où $T$ est l'ensemble des éléments
\emph{homogènes} de $S_+$ qui n'appartiennent pas à $\mathfrak{p}$.
\end{env}

\subsection{Spectre premier homogène d'un anneau gradué.}
\label{subsection:II.2.3-fr}

\begin{env}[2.3.1]
\label{II.2.3.1-fr}
Étant donné un anneau gradué $S$, à degrés positifs, on appelle
\emph{spectre premier homogène} de $S$ et on désigne par
$\operatorname{Proj}(S)$ l'ensemble des idéaux premiers gradués de $S_+$
(\hyperref[II.2.1.10-fr]{2.1.10}), ou, ce qui revient au même, l'ensemble
des idéaux premiers gradués de $S$ \emph{ne contenant pas} $S_+$~; nous
allons définir une structure de \emph{schéma} ayant
$\operatorname{Proj}(S)$ comme ensemble de base.
\end{env}

\begin{env}[2.3.2]
\label{II.2.3.2-fr}
Pour toute partie $E$ de $S$, soit $V_+(E)$ l'ensemble des idéaux premiers
gradués de $S$ contenant $E$ et ne contenant pas $S_+$~; c'est donc la
partie $V(E)\cap\operatorname{Proj}(S)$ de
$\operatorname{Spec}(S)$. De (\hyperref[I.1.1.2-fr]{I, 1.1.2}) on déduit
donc~:
\[
  \label{II.2.3.2.1-fr}
  V_+(0)=\operatorname{Proj}(S),\quad
  V_+(S)=V_+(S_+)=\varnothing
  \tag{2.3.2.1}
\]
\[
  \label{II.2.3.2.2-fr}
  V_+\left(\bigcup_\lambda E_\lambda\right)
  =\bigcap_\lambda V_+(E_\lambda)
  \tag{2.3.2.2}
\]
\[
  \label{II.2.3.2.3-fr}
  V_+(EE')=V_+(E)\cup V_+(E')
  \tag{2.3.2.3}
\]
On ne change pas $V_+(E)$ en remplaçant $E$ par l'idéal gradué engendré par
$E$~; en outre, si $\mathfrak{J}$ est un idéal gradué de $S$, on a
\[
  \label{II.2.3.2.4-fr}
  V_+(\mathfrak{J})
  =V_+\left(\bigcup_{q\geq n}(\mathfrak{J}\cap S_q)\right)
  \tag{2.3.2.4}
\]
